% =====================================================================
%  I. Introduction
%
%  Ban da nop mo bai bang "QSVM dat 0.854, hon SVM-RBF 0.838". So do gio
%  khong con dung (XGBoost 0.8503 dung dau). Ban nay mo bang CAU HOI, va
%  noi thang ngay o trang 1 rang phan lon ket qua la am hoac khong ket luan
%  duoc. Reviewer doc trang 1 phai thay ngay minh khong ban loi the.
%
%  Dong cuoi cua danh sach dong gop la cho tu khai 5 khang dinh phai rut.
%  De o day chu khong giau xuong muc VII: neu reviewer tu tim ra thi mat
%  het thien chi.
% =====================================================================

\revnote{rewrite}{Rewritten. The submitted introduction claimed a low-data and prior-shift advantage; both are withdrawn (AE-1, R1-4).}
\begin{revbody}
Network intrusion detection is an attractive target for quantum machine
learning. The data are tabular and moderate-dimensional, the decision is a
classification, and the operational value of a small accuracy gain is high
\cite{ref24}. Classical machine learning and deep learning are well established
on the task \cite{ref1,ref2,ref3}, which sets the bar a quantum method has to
clear. Quantum kernel methods are the natural instrument: a feature map prepares a
state whose fidelities form a positive semi-definite kernel that can be dropped
into a standard support vector machine \cite{havlicek2019,schuld2019}, so the
quantum component is a single well-isolated block inside an otherwise classical
pipeline.

The literature reflects that attractiveness. Studies applying quantum kernels
and variational classifiers to NSL-KDD \cite{ref4} and its successors
\cite{ref25} report accuracies competitive with, and sometimes above,
classical baselines \cite{ref13,ref14,ref15}. What they rarely report is the shape of the claim: at which training
set size, against which class prior, at which circuit width, and against which
classical opponent. Those choices are not incidental. We show below that varying
only the training-set size, with everything else fixed, is enough to reverse the
ordering between a quantum kernel and the tree ensembles and the linear and
polynomial kernels we tested, on a shared four-dimensional embedding.
A benchmark conducted at one operating point can therefore report an advantage
or a deficit and be internally consistent in either case.

This paper asks a narrower question than the literature usually asks:
\emph{where} is a NISQ-feasible quantum kernel worth its cost? The question is
posed as a regime-identification problem (Problem~\ref{prob:regime}) whose
answers include the regimes where the answer is ``nowhere we can detect''.

Our protocol varies one factor at a time. The embedding dimension is fixed by a
lexicographic rule that charges explicitly for two-qubit gates, before any
downstream metric is inspected. The entangling layer is ablated against an
otherwise identical control. The test distribution is shifted four ways. The
training-set size is swept over two orders of magnitude under both the natural
class prior and a rare-enriched one. Every comparison is paired within run over
ten runs, both quantum and classical models are tuned by the same procedure on
the same held-out set, and $p$-values are Holm-corrected within baseline
families. The baselines include a random forest and XGBoost, which the submitted
version of this work did not.

The headline is a map, not a number. Of $110$ controlled comparisons, $21$
favour the quantum kernel, $21$ favour a classical baseline, and $68$ are
inconclusive. Inside that map there is structure worth reporting:

\begin{itemize}
  \item \textbf{An ordering that reverses with sample size.} Under the natural
    class prior, on the shared four-component embedding, \QSVM{} is the
    second-worst model at $N=100$ and the best at $N=10^{4}$; the paired
    difference against XGBoost and the random forest changes sign between
    $N=2000$ and $N=5000$ in both tuning arms and reaches corrected
    significance, the sign change against SVM-RBF is $+0.0008$ and never
    significant, and against XGBoost given all $122$ features the kernel
    reaches parity rather than a lead. Whether rare-attack enrichment (the
    sampling our own submitted protocol used) removes the reversal is not
    decidable with this design: the enriched arm ends at $N=2000$, before the
    reversal, and where the two arms overlap the tree ensembles lead under
    both priors.
  \item \textbf{A selection rule that transfers.} The three-stage rule returns
    $n^{\ast}=4$ on NSL-KDD and, applied without changing a threshold,
    $n^{\ast}=6$ on UNSW-NB15, recovered on $10/10$ independent subsets. The
    entanglement ablation transfers as well; the sample-size reversal does not.
  \item \textbf{A measured boundary.} Spending a larger feature budget forces a
    wider circuit, and widening destroys the kernel: at $K=80$, where the rule
    returns $n^{\ast}=8$, all $48$ paired comparisons favour the classical
    model and the entanglement ablation reverses sign. The mechanism is visible
    in the Gram matrix, whose off-diagonal dispersion decays faster for the
    \ZZ{} map than for its entanglement-free control: by a factor of $2.02$ at
    $K=20$, the ratio the $\binom{n}{2}$ versus $n$ phase-term counting
    predicts, and by $3.32$ at $K=80$, where the collapse occurs and the
    counting argument predicts the ordering but not the rate. Across widths the
    dispersion correlates with macro-$F_{1}$ at $r=+0.77$ for the \ZZ{} map
    (seven widths, 95\% CI $[+0.04,+0.96]$); the contrast with the \Zmap{}
    map ($r=+0.32$) is not resolvable at seven points.
\end{itemize}

We also report where the method loses. It degrades faster than either tree
ensemble under class-prior shift; it is the least robust of all seven models to
feature perturbation; and although it detects rare attacks far better than the
tree ensembles, an RBF-kernel SVM detects them better still.
\end{revbody}

\subsection*{Relation to the submitted version}
\revnote{new}{New subsection: the five claims of the submitted version that are withdrawn, stated on page 1.}
\begin{revbody}

This paper is a substantially revised version of an earlier submission, and five
of that version's claims do not survive the revised protocol. We state them
here rather than in a footnote. (i)~The claimed advantage in the low-data regime
is reversed in direction. (ii)~A reported rare-attack gain of $6.7$ points with
$d=+0.68$ is withdrawn: no rare-class metric was computed by the code that
produced it. (iii)~The former Theorem~1 was false, as one reviewer identified;
it and the objective it maximised have been removed rather than repaired.
(iv)~The former Pareto stage filtered no candidate on this data and has been
replaced. (v)~With a tree ensemble in the comparison, \QSVM{} is no longer the
top-ranked model at the main operating point. Sec.~\ref{sec:erratum} gives the
technical detail, including one error no reviewer raised.
\end{revbody}
